% Section 8: drafting (drafter workstream draft, integrated by the paper agent).
% Claims and status:
%  - Checkpoint identity, architecture, conditioning layers, mask token: config and weights at
%    the pinned revision (finding; evidence/drafter/README.md pending commit).
%  - How SGLang runs a DFlash cycle on the GDN hybrid: code reading at bd66ce34.
%  - Parameter and byte counts per cycle: derived from the configs.
%  - Card numbers (B200): the authors' claims.
%  - All GH200 measurements: pending (drafter).
\section{Drafting}\label{sec:drafting}
How much a verify pass is worth depends on the drafter in front of it: a cycle commits as many tokens as the target accepts plus one, and its cost includes the draft. The native MTP layer of Qwen3.5-4B drafts sequentially, paying one small layer plus the full 1.27~GB head per drafted token, and its acceptance decays with depth. A block drafter proposes a whole block in one forward pass. This section describes the public block drafter for this target, how SGLang runs it on a hybrid model, and what is measured about it.

\subsection{The public block drafter}\label{sec:drafter-model}
DFlash has a public drafter for exactly this target, \code{z-lab/Qwen3.5-4B-DFlash} at revision \texttt{9a1996cc}, under Apache-2.0~\cite{dflashpaper,dflash4bcard}; \code{modal-labs/Qwen3.5-4B-DFlash} carries byte-identical weights. Read from the checkpoint rather than the card, it is a six-layer Qwen3-style transformer of width 2,560 (32 query and 8 key--value heads of dimension 128, MLP width 9,216) with no embedding or head of its own: it embeds with the target's table and projects through the target's tied head. It conditions on the target's hidden states after layers 1, 5, 9, 13, 17, 21, 25 and 29, concatenated per token and projected by one linear map and an RMS norm to a single context vector. Five layers use sliding-window attention (window 4,096) and the last uses full attention. It has 634 million parameters, 1.27~GB in BF16, almost exactly the size of the target's head (derived).

A cycle in SGLang's DFlash worker~\cite{sgcode} works as follows (code reading at the pin). The block is the last committed token followed by fifteen copies of a mask token (id 248,077). Each draft layer's keys and values for the committed context come from the projected target features, which the worker writes into the drafter's own cache as tokens are committed. One forward pass over the sixteen positions yields sixteen hidden states; positions 1--15 go through the target head, and their argmaxes are the drafted tokens. The target verifies the block in one pass, the worker accepts the longest drafted prefix that matches the target's argmax, commits it plus the target's token at the first mismatch, and appends the committed tokens' features to the drafter's cache. The sliding layers are causal inside the block and the full layer is not; SGLang infers this from the layer types because the configuration sets no \code{is\_causal}, which matches the attention pattern of the SpecForge training recipe. % CITE: SpecForge (sgl-project/SpecForge, commit 3cb0510), DFlash training masks

On the hybrid target the verify pass has a cost the drafter's parameters do not show. SGLang verifies a chain on the GDN layers by keeping the recurrent state after every block position, so that the state after the accepted prefix can be committed without recomputation. Each state is 50.3~MB per request, so a 16-token block writes about 805~MB per request per cycle, and each running request reserves space for sixteen of them (derived). That bounds block-16 capacity on this machine, and it is a property of the verifier, not of the drafter.

The card's launch line targets a B200 with TRT-LLM target attention, FA4 draft attention and piecewise-compiled prefill graphs. None carries over unchanged: TRT-LLM attention is Blackwell-only, FA3 and FA4 are unavailable on this aarch64 host, and the piecewise prefill graph fails to capture for this target at the pinned commit. The drafter is served with FlashInfer for target and draft attention, the FlashInfer GDN kernels, the default prefill graph and the card's overlap planning stream, with decode, verify and draft CUDA graphs and the overlap scheduler on \pending{drafter: launch record and graph-capture state for every run; \nolinkurl{evidence/drafter/launch/}}.

\subsection{What a cycle costs}\label{sec:drafter-cost}
Counting weight bytes gives the shape of the trade before any timing. At batch one a DFlash cycle reads the drafter once (1.27~GB), the target head once for the fifteen drafted positions (1.27~GB), and the target once to verify (8.41~GB): drafting costs about 30\% of the verify pass's weight traffic and proposes fifteen tokens. An MTP chain of three steps reads the MTP layer and the full head three times (about 4.5~GB, 54\% of the verify pass) and proposes three (derived). Activation traffic, the per-position GDN states, attention over the context and host time are not in these counts \pending{drafter: split of the DFlash cycle at batch 1, 4, 16 and 64 into draft forward, draft head, verification, GDN commit, cache update and host time; \nolinkurl{evidence/drafter/cycle_split.csv}}.

\subsection{Acceptance by block position}\label{sec:drafter-acceptance}
Let $h_j$ be the number of verify cycles in which exactly $j$ drafted tokens were accepted, $N=\sum_jh_j$, and $S(k)=\sum_{j\ge k}h_j/N$ the fraction of cycles that accept at least $k$. The conditional acceptance of block position $k$ is $\alpha_k=S(k)/S(k-1)$, and a cycle commits $\tau=1+\sum_{k\ge1}S(k)$ tokens on average. SGLang reports $h$ per request, so $\alpha_k$ and $\tau$ are measured rather than modelled. The card reports mean accepted lengths of 5.9--7.7 at block 16 on a B200 for its five benchmarks, against 3.3--3.5 for MTP with three steps~\cite{dflash4bcard}; these are the authors' numbers on other hardware.

\pending{drafter: $\alpha_k$ for $k=1,\dots,15$ at block sizes 4, 8 and 16 per domain, greedy with thinking on, concurrency 1, and $\tau$ against tuned MTP on the same prompts; \nolinkurl{evidence/drafter/acceptance_by_position.csv}}

\pendingfigure{Conditional acceptance $\alpha_k$ against block position $k$, one line per domain, block size 16, with MTP's three positions for comparison.}{evidence/drafter/acceptance_by_position.csv}

\subsection{Output equality and limits}\label{sec:drafter-limits}
Greedy speculative verification is exact in real arithmetic, but the verify pass computes the target's logits at a different batch shape from plain decoding, so the stock noise floor of Section~\ref{sec:divergence} applies \pending{drafter: divergences of DFlash (blocks 8 and 16) against plain decoding per 1,000 tokens, classified by mechanism, against the plain-decoding floor; \nolinkurl{evidence/drafter/equality.json}}. Four limitations are candidates: acceptance that decays with block position or by domain; drafting cost dominated by the full-vocabulary head for fifteen positions; the verifier's per-position state writes, which bound block size at high concurrency; and draft states far from the target's in the head's metric, which decides whether transport can use DFlash drafts (Section~\ref{sec:transport-measured}). A custom drafter is worth training only against one of these, judged by total serving time: fine-tuning on the target's own outputs for acceptance, adding DFlash~2's selector~\cite{dflash} for acceptance, or an auxiliary loss on $\norm{h^t-h^d}_2^2/\norm{h^t}_2^2$ if transport's required drift reduction proves reachable \pending{drafter: which limitation the measurements support, and the decision on training}.
